% =====================================================================
% PM QLKT — Slides bảo vệ ĐATN (Beamer + xelatex)
% =====================================================================
% Build: xelatex defense.tex (chạy 2 lần để cập nhật cross-ref + TOC)
% =====================================================================

\documentclass[aspectratio=43,xcolor=table]{beamer}

% --- Font + tiếng Việt ---
\usepackage{fontspec}
\usepackage{xcolor}
\usepackage{tikz}
\usetikzlibrary{positioning,arrows.meta,shapes.geometric}
\usepackage{array}
\usepackage{booktabs}
\usepackage{tabularx}
\usepackage{listings}

% Vietnamese support
\usepackage{polyglossia}
\setdefaultlanguage{vietnamese}
\setotherlanguage{english}
\setmainfont{Arial}
\setsansfont{Arial}
\setmonofont{Menlo}

% --- HUST color palette ---
\definecolor{hustred}{RGB}{200,16,46}
\definecolor{hustdark}{RGB}{138,12,32}
\definecolor{hustlight}{RGB}{255,232,236}
\definecolor{bgsoft}{RGB}{250,250,250}
\definecolor{linecol}{RGB}{229,231,235}
\definecolor{textmain}{RGB}{31,41,55}
\definecolor{textmuted}{RGB}{107,114,128}

% --- Beamer theme + tuỳ biến ---
\usetheme{default}
\usecolortheme[named=hustred]{structure}
\useinnertheme{rectangles}
\useoutertheme{infolines}

% Tuỳ biến header/footer
\setbeamercolor{title}{fg=hustred}
\setbeamercolor{frametitle}{fg=white,bg=hustred}
\setbeamercolor{section in head/foot}{fg=white,bg=hustdark}
\setbeamercolor{subsection in head/foot}{fg=white,bg=hustred}
\setbeamercolor{author in head/foot}{fg=white,bg=hustdark}
\setbeamercolor{title in head/foot}{fg=white,bg=hustred}
\setbeamercolor{date in head/foot}{fg=white,bg=hustdark}

\setbeamercolor{block title}{fg=white,bg=hustred}
\setbeamercolor{block body}{fg=textmain,bg=hustlight}
\setbeamercolor{block title alerted}{fg=white,bg=hustdark}
\setbeamercolor{block body alerted}{fg=textmain,bg=bgsoft}
\setbeamercolor{block title example}{fg=white,bg=hustdark}
\setbeamercolor{block body example}{fg=textmain,bg=bgsoft}

\setbeamercolor{itemize item}{fg=hustred}
\setbeamercolor{itemize subitem}{fg=hustdark}
\setbeamertemplate{itemize item}[circle]
\setbeamertemplate{itemize subitem}[triangle]

% Frame title style
\setbeamertemplate{frametitle}{%
  \nointerlineskip%
  \begin{beamercolorbox}[wd=\paperwidth,ht=2.5ex,dp=1ex,leftskip=0.3cm,rightskip=0.3cm]{frametitle}%
    \strut\insertframetitle\strut%
  \end{beamercolorbox}%
}

\setbeamertemplate{navigation symbols}{}

% Caption "X / Y" thay vì page number
\setbeamertemplate{footline}{%
  \leavevmode%
  \hbox{%
    \begin{beamercolorbox}[wd=.4\paperwidth,ht=2.25ex,dp=1ex,center]{author in head/foot}%
      \usebeamerfont{author in head/foot}\insertshortauthor%
    \end{beamercolorbox}%
    \begin{beamercolorbox}[wd=.45\paperwidth,ht=2.25ex,dp=1ex,center]{title in head/foot}%
      \usebeamerfont{title in head/foot}PM QLKT --- ĐATN 2025/2026%
    \end{beamercolorbox}%
    \begin{beamercolorbox}[wd=.15\paperwidth,ht=2.25ex,dp=1ex,center]{date in head/foot}%
      \usebeamerfont{date in head/foot}\insertframenumber{} / \inserttotalframenumber%
    \end{beamercolorbox}}%
}

% --- Macros tiện dụng ---
\newcommand{\highlight}[1]{\textcolor{hustred}{\textbf{#1}}}
\newcommand{\hustpill}[1]{\colorbox{hustlight}{\textcolor{hustred}{\small #1}}}

% Listings setup (nếu có code)
\lstdefinestyle{tscode}{
  basicstyle=\ttfamily\scriptsize\color{textmain},
  keywordstyle=\color{hustred}\bfseries,
  commentstyle=\color{textmuted},
  stringstyle=\color{hustdark},
  backgroundcolor=\color{bgsoft},
  frame=single,
  rulecolor=\color{linecol},
  showstringspaces=false,
  breaklines=true,
  tabsize=2,
}
\lstset{style=tscode}

% --- Metadata ---
\title{Phần mềm Quản lý Khen thưởng}
\subtitle{Tại Học viện Khoa học Quân sự}
\author[Trần Anh Đức]{Sinh viên thực hiện: \textbf{Trần Anh Đức}\\Giảng viên hướng dẫn: \textbf{[Học hàm. Tên giảng viên]}}
\institute[ĐHBKHN]{Trường Công nghệ Thông tin và Truyền thông --- Đại học Bách khoa Hà Nội}
\date{Tháng 5/2026}

% =====================================================================
\begin{document}

% --- Slide 1: Title ---
{
\setbeamertemplate{footline}{}
\begin{frame}[plain]
  \begin{center}
    \vspace{0.5cm}
    {\large \textcolor{hustdark}{ĐẠI HỌC BÁCH KHOA HÀ NỘI}}\\
    {\small \textcolor{textmuted}{Trường Công nghệ Thông tin và Truyền thông}}\\[1.2cm]

    {\Huge \textcolor{hustred}{\textbf{PHẦN MỀM QUẢN LÝ}}}\\[0.3cm]
    {\Huge \textcolor{hustred}{\textbf{KHEN THƯỞNG}}}\\[0.5cm]

    {\large Đồ án Tốt nghiệp Đại học --- Niên khóa 2025/2026}\\[1.5cm]

    \begin{tabular}{rl}
      \textbf{Sinh viên thực hiện} & : Trần Anh Đức \\
      \textbf{Mã số sinh viên} & : [MSSV] \\
      \textbf{Lớp} & : [Lớp] \\
      \textbf{Giảng viên hướng dẫn} & : [Học hàm. Tên GVHD] \\
    \end{tabular}\\[1cm]

    {\small \textcolor{textmuted}{Hà Nội, tháng 5 năm 2026}}
  \end{center}
\end{frame}
}

% --- Slide 2: Nội dung trình bày ---
\begin{frame}{Nội dung trình bày}
  \begin{columns}[T,onlytextwidth]
    \column{0.5\textwidth}
    \textcolor{hustdark}{\textbf{Phần I --- Mở đầu}}
    \begin{enumerate}
      \item Mục tiêu của đồ án
      \item Phân tích bài toán
      \item Hệ thống danh hiệu khen thưởng
      \item Logic chuỗi danh hiệu
    \end{enumerate}
    \vspace{0.3cm}
    \textcolor{hustdark}{\textbf{Phần II --- Thiết kế}}
    \begin{enumerate}\setcounter{enumi}{4}
      \item Kiến trúc \& công nghệ
      \item Cơ sở dữ liệu
      \item Phân tích chức năng \& quy trình
      \item Kiến trúc phần mềm
    \end{enumerate}

    \column{0.5\textwidth}
    \textcolor{hustdark}{\textbf{Phần III --- Tính năng nổi bật}}
    \begin{enumerate}\setcounter{enumi}{8}
      \item Tự động xét điều kiện
      \item Quy trình đề xuất \& ra quyết định
      \item Nhập/xuất dữ liệu hàng loạt
      \item Vận hành \& quản trị hệ thống
    \end{enumerate}
    \vspace{0.3cm}
    \textcolor{hustdark}{\textbf{Phần IV --- Đánh giá \& Tổng kết}}
    \begin{enumerate}\setcounter{enumi}{12}
      \item Kiểm thử, hiệu năng, triển khai
      \item Kết luận và hướng phát triển
    \end{enumerate}
  \end{columns}

  \vspace{0.4cm}
  \begin{block}{}
    \centering Thời lượng trình bày: \highlight{13--14 phút}, sau đó là phần Hỏi \& Đáp với Hội đồng.
  \end{block}
\end{frame}

% --- Slide 3: Mục tiêu của đồ án ---
\begin{frame}{Mục tiêu của đồ án}
  \begin{columns}[T,onlytextwidth]
    \column{0.5\textwidth}
    \textbf{Bối cảnh và động lực}
    \begin{itemize}\small
      \item Quản lý khen thưởng tại đơn vị quân đội hiện chủ yếu dựa trên \textbf{Excel và hồ sơ giấy}
      \item Khó tra cứu lịch sử khen thưởng nhiều năm
      \item \highlight{Quy định chuỗi danh hiệu phức tạp} (BKBQP, CSTĐTQ, BKTTCP) dễ sai sót khi tính thủ công
      \item Vận hành trên \textbf{mạng nội bộ}, không phụ thuộc Internet
    \end{itemize}

    \column{0.5\textwidth}
    \textbf{Mục tiêu cụ thể}
    \begin{itemize}\small
      \item Quản lý \textbf{5 nhóm khen thưởng} (UC5--UC9) với 4 cấp vai trò
      \item \textbf{Tự động hoá} xét điều kiện theo quy định hiện hành
      \item Hỗ trợ \textbf{đề xuất $\to$ duyệt $\to$ ra quyết định}
      \item Đảm bảo \textbf{truy vết và bảo mật} đầy đủ
      \item \textbf{Triển khai trên mạng nội bộ} offline
    \end{itemize}
    \begin{block}{Mục tiêu cốt lõi}
      \small Giảm sai sót thủ công, tiết kiệm thời gian, đảm bảo đúng quy định.
    \end{block}
  \end{columns}
\end{frame}

% --- Slide 4: Phân tích bài toán — 4 thách thức cốt lõi ---
\begin{frame}{Phân tích bài toán --- 4 thách thức cốt lõi}
  \begin{columns}[T,onlytextwidth]
    \column{0.5\textwidth}
    \textbf{\textcolor{hustred}{\textcircled{1} Logic chuỗi danh hiệu nhiều năm}}\\
    \scriptsize Cửa sổ trượt 3y/7y, ràng buộc lifetime, lỡ đợt --- không thể xét bằng truy vấn đơn lẻ, đòi hỏi \textbf{engine xét điều kiện trên toàn timeline} của quân nhân.\\[0.3cm]

    \textbf{\textcolor{hustred}{\textcircled{2} Tính nhất quán đa bảng}}\\
    \scriptsize Một lần duyệt đề xuất cập nhật $\geq$ 5 bảng (đề xuất, hồ sơ, số quyết định, file PDF, thông báo, audit). Bắt buộc \textbf{transaction} và \textbf{chống race condition} khi hai cán bộ duyệt cùng lúc.

    \column{0.5\textwidth}
    \textbf{\textcolor{hustred}{\textcircled{3} Phân quyền theo cây tổ chức}}\\
    \scriptsize 4 vai trò $\times$ cây đơn vị nhiều cấp $\to$ \textbf{hai lớp middleware độc lập} (\texttt{requireRole} + \texttt{unitFilter}). Bypass được lớp này thì lớp kia vẫn chặn.\\[0.3cm]

    \textbf{\textcolor{hustred}{\textcircled{4} Vận hành mạng nội bộ (offline-first)}}\\
    \scriptsize Không CDN, không dịch vụ ngoài $\to$ tự host font, tự backup theo lịch, \textbf{PM2 auto-restart}, cài đặt qua thư mục \texttt{node\_modules} đóng gói sẵn.
  \end{columns}

  \vspace{0.4cm}
  \begin{block}{}
    \centering\small Trọng tâm trình bày: \highlight{\textcircled{1}} (slide 6--7) và \highlight{\textcircled{2}} (slide 13, 18).
  \end{block}
\end{frame}

% --- Slide 5: Hệ thống khen thưởng — 5 nhóm UC5–UC9 ---
\begin{frame}{Hệ thống khen thưởng --- 5 nhóm UC5--UC9}
  \scriptsize
  \begin{tabularx}{\textwidth}{|c|>{\raggedright\arraybackslash}p{2.4cm}|>{\raggedright\arraybackslash}X|>{\raggedright\arraybackslash}p{2.6cm}|}
    \hline
    \rowcolor{hustred}\textcolor{white}{\textbf{UC}} & \textcolor{white}{\textbf{Nhóm khen thưởng}} & \textcolor{white}{\textbf{Phân rã chi tiết}} & \textcolor{white}{\textbf{Tiêu chí xét}} \\ \hline
    \textbf{UC5} & \textbf{Khen thưởng hằng năm} & Cá nhân (CSTĐCS, CSTT, BKBQP, CSTĐTQ, BKTTCP) + Đơn vị (ĐVQT, ĐVTT, BKBQP, BKTTCP) & Theo năm \\ \hline
    \textbf{UC6} & \textbf{Khen thưởng niên hạn} & HCCSVV (3 hạng), HCQKQT, KNC VSNXD QĐNDVN & Theo thời gian phục vụ \\ \hline
    \textbf{UC7} & \textbf{Khen thưởng cống hiến} & HCBVTQ (3 hạng) & Cộng dồn 120 tháng \\ \hline
    \textbf{UC8} & \textbf{Khen thưởng thành tích} & NCKH (Đề tài KH, Sáng kiến KH) & Theo công trình \\ \hline
    \textbf{UC9} & \textbf{Khen thưởng đột xuất} & Theo sự kiện / chiến công cụ thể & Phi định kỳ \\ \hline
  \end{tabularx}

  \vspace{0.4cm}
  \begin{block}{}
    Mỗi nhóm UC được hiện thực bằng một hoặc nhiều \highlight{proposal type} trong code, dispatch qua \highlight{Strategy pattern} (Slide 13).
  \end{block}
\end{frame}

% --- Slide 6: Khen thưởng hằng năm — cá nhân + đơn vị (timeline) ⭐ ---
\begin{frame}{Khen thưởng hằng năm --- chu trình hoàn hảo}
  \centering

  % ===== CÁ NHÂN lane =====
  \begin{flushleft}\hspace{0.2cm}\textbf{\textcolor{hustdark}{Cá nhân}} \textit{\scriptsize(CSTĐCS + NCKH liên tục mỗi năm là nền tảng; đề xuất chuỗi xét dựa trên streak năm trước)}\end{flushleft}
  \vspace{-0.15cm}
  \begin{tikzpicture}[
    every node/.style={font=\scriptsize},
    base/.style={draw=hustred!60, fill=hustred!5, rounded corners=1pt, font=\tiny, inner sep=1.5pt, align=center},
    award/.style={draw=hustred, fill=hustred!10, rounded corners=2pt, font=\tiny, inner sep=2pt, align=center},
    awardhigh/.style={draw=hustdark, fill=hustred!25, rounded corners=2pt, font=\tiny\bfseries, inner sep=2pt, align=center},
    awardtop/.style={draw=hustdark, fill=hustdark, text=white, rounded corners=3pt, font=\tiny\bfseries, inner sep=2pt, align=center},
    yearmark/.style={draw=hustdark, thick}
  ]
    \draw[->, very thick, hustdark] (-0.3,0) -- (10.0,0);
    \foreach \i in {1,...,8} {
      \pgfmathsetmacro\xcoord{(\i-1)*1.15}
      \draw[yearmark] (\xcoord, 0.08) -- (\xcoord, -0.08);
      \node[above=0.05cm, font=\tiny\bfseries] at (\xcoord, 0.08) {Năm \i};
      \node[base, anchor=north] at (\xcoord, -0.18) {CSTĐCS};
    }
    \node[right, font=\tiny] at (10.0, -0.05) {\dots};

    % Năm 3: BKBQP#1
    \node[award, anchor=south] (b1) at (2.30, 0.65) {BKBQP \#1};
    \draw[hustred, dashed, thin] (2.30, 0.5) -- (b1.south);

    % Năm 4: CSTĐTQ#1
    \node[awardhigh, anchor=south] (c1) at (3.45, 1.3) {CSTĐTQ \#1};
    \draw[hustdark, dashed, thin] (3.45, 0.5) -- (c1.south);

    % Năm 5: BKBQP#2
    \node[award, anchor=south] (b2) at (4.60, 0.65) {BKBQP \#2};
    \draw[hustred, dashed, thin] (4.60, 0.5) -- (b2.south);

    % Năm 7: BKBQP#3 + CSTĐTQ#2
    \node[awardhigh, anchor=south] (c2) at (6.90, 1.3) {BKBQP \#3 + CSTĐTQ \#2};
    \draw[hustdark, dashed, thin] (6.90, 0.5) -- (c2.south);

    % Năm 8: BKTTCP
    \node[awardtop, anchor=south] (bk) at (8.05, 2.0) {$\bigstar$ BKTTCP (lifetime)};
    \draw[hustdark, very thick] (8.05, 0.5) -- (bk.south);
  \end{tikzpicture}

  \vspace{0.25cm}

  % ===== ĐƠN VỊ lane =====
  \begin{flushleft}\hspace{0.2cm}\textbf{\textcolor{hustdark}{Đơn vị}} \textit{\scriptsize(ĐVQT liên tục; không CSTĐTQ, không bắt buộc NCKH, BKTTCP lặp lại)}\end{flushleft}
  \vspace{-0.15cm}
  \begin{tikzpicture}[
    every node/.style={font=\scriptsize},
    base/.style={draw=hustred!60, fill=hustred!5, rounded corners=1pt, font=\tiny, inner sep=1.5pt, align=center},
    uaward/.style={draw=hustred, fill=hustred!10, rounded corners=2pt, font=\tiny, inner sep=2pt, align=center},
    uawardtop/.style={draw=hustdark, fill=hustred!40, rounded corners=3pt, font=\tiny\bfseries, inner sep=2pt, align=center},
    yearmark/.style={draw=hustdark, thick}
  ]
    \draw[->, very thick, hustdark] (-0.3,0) -- (10.0,0);
    \foreach \i in {1,...,8} {
      \pgfmathsetmacro\xcoord{(\i-1)*1.15}
      \draw[yearmark] (\xcoord, 0.08) -- (\xcoord, -0.08);
      \node[above=0.05cm, font=\tiny\bfseries] at (\xcoord, 0.08) {Năm \i};
      \node[base, anchor=north] at (\xcoord, -0.18) {ĐVQT};
    }
    \node[right, font=\tiny] at (10.0, -0.05) {\dots\ \textit{(lặp lại)}};

    % Năm 3, 5, 7: BKBQP đv
    \node[uaward, anchor=south] (ub1) at (2.30, 0.65) {BKBQP đv \#1};
    \draw[hustred, dashed, thin] (2.30, 0.5) -- (ub1.south);

    \node[uaward, anchor=south] (ub2) at (4.60, 0.65) {BKBQP đv \#2};
    \draw[hustred, dashed, thin] (4.60, 0.5) -- (ub2.south);

    \node[uaward, anchor=south] (ub3) at (6.90, 0.65) {BKBQP đv \#3};
    \draw[hustred, dashed, thin] (6.90, 0.5) -- (ub3.south);

    % Năm 8: BKTTCP đv
    \node[uawardtop, anchor=south] (ubk) at (8.05, 1.3) {BKTTCP đv \\ (lặp lại)};
    \draw[hustdark, very thick] (8.05, 0.5) -- (ubk.south);
  \end{tikzpicture}

  \vspace{0.1cm}
  \scriptsize\centering
  \textcolor{hustdark}{\textbf{Lỡ đợt:}} chưa đề xuất ở năm đến mốc (vd Năm 2) $\to$ chu kỳ vẫn đếm tiếp $\to$ Năm 4 lại được xét, không phải khôi phục chuỗi từ đầu.
\end{frame}

% --- Slide 7: Các loại khen thưởng còn lại — niên hạn + 1-lần + cống hiến ⭐ ---
\begin{frame}{Bốn loại khen thưởng dài hạn}
  \begin{columns}[T,onlytextwidth]
    \column{0.58\textwidth}

    \textbf{\textcolor{hustdark}{Niên hạn + 1-lần (theo năm phục vụ)}}\\[0.05cm]
    \begin{tikzpicture}[
      every node/.style={font=\scriptsize},
      nh/.style={draw=hustred, fill=hustred!10, rounded corners=2pt, font=\tiny, inner sep=1.5pt, align=center, minimum width=1.1cm},
      once/.style={draw=hustdark, fill=hustdark, text=white, rounded corners=3pt, font=\tiny\bfseries, inner sep=1.5pt, align=center, minimum width=1.1cm},
      knc/.style={draw=hustdark, fill=hustred!40, rounded corners=3pt, font=\tiny\bfseries, inner sep=1.5pt, align=center, minimum width=1.1cm},
      yearmark/.style={draw=hustdark, thick}
    ]
      \draw[->, very thick, hustdark] (0,0) -- (6.0,0);
      \foreach \i [count=\j from 0] in {10, 15, 20, 25} {
        \pgfmathsetmacro\xcoord{0.4 + \j*1.55}
        \draw[yearmark] (\xcoord, 0.1) -- (\xcoord, -0.1);
        \node[below=0.06cm, font=\tiny\bfseries] at (\xcoord, 0) {Năm \i};
      }

      % Year 10: HCCSVV hạng Ba
      \node[nh, anchor=south] (n1) at (0.4, 0.3) {HCCSVV \\ Ba};
      % Year 15: HCCSVV hạng Nhì
      \node[nh, anchor=south] (n2) at (1.95, 0.3) {HCCSVV \\ Nhì};
      % Year 20: HCCSVV hạng Nhất + KNC nữ
      \node[nh, anchor=south] (n3) at (3.5, 0.3) {HCCSVV \\ Nhất};
      \node[knc, above=0.04cm of n3] (kn) {KNC (nữ)};
      % Year 25: HCQKQT + KNC nam
      \node[knc, anchor=south] (km) at (5.05, 0.3) {KNC (nam)};
      \node[once, above=0.04cm of km] (q) {HCQKQT};

      \draw[hustred, dashed, thin] (0.4, 0.1) -- (n1.south);
      \draw[hustred, dashed, thin] (1.95, 0.1) -- (n2.south);
      \draw[hustred, dashed, thin] (3.5, 0.1) -- (n3.south);
      \draw[hustdark, very thick] (5.05, 0.1) -- (km.south);
    \end{tikzpicture}

    \vspace{0.15cm}
    \scriptsize
    \textbf{HCCSVV} 3 hạng theo mốc 10/15/20y. \textbf{HCQKQT} 25y --- 1 lần.\\
    \textbf{KNC VSNXD QĐNDVN}: nam 25y / nữ 20y --- 1 lần duy nhất.

    \column{0.42\textwidth}

    \textbf{\textcolor{hustdark}{Cống hiến (HCBVTQ)}}\\[0.05cm]
    \scriptsize Tích luỹ \textbf{120 tháng} giữ chức vụ --- không phụ thuộc năm phục vụ, mà theo \emph{nhóm hệ số chức vụ}:\\[0.15cm]

    \begin{tikzpicture}[
      every node/.style={font=\tiny},
      grpbox/.style={draw=hustred, fill=hustred!10, rounded corners=2pt, inner sep=3pt, minimum width=1.6cm, font=\tiny\bfseries, align=center},
      tierbox/.style={draw=hustdark, fill=hustred!30, rounded corners=2pt, inner sep=3pt, minimum width=1.4cm, font=\tiny\bfseries, align=center}
    ]
      % Row 1: nhóm 0.7 → Hạng Ba
      \node[grpbox] (g1) at (0,0) {Nhóm 0.7};
      \node[tierbox] (t1) at (3.4,0) {Hạng Ba};
      \draw[->, thick, hustdark] (g1.east) -- node[above, font=\tiny]{120 tháng} (t1.west);

      % Row 2: nhóm 0.8 → Hạng Nhì
      \node[grpbox] (g2) at (0,-0.9) {Nhóm 0.8};
      \node[tierbox] (t2) at (3.4,-0.9) {Hạng Nhì};
      \draw[->, thick, hustdark] (g2.east) -- node[above, font=\tiny]{120 tháng} (t2.west);

      % Row 3: nhóm 0.9-1.0 → Hạng Nhất
      \node[grpbox] (g3) at (0,-1.8) {Nhóm 0.9--1.0};
      \node[tierbox] (t3) at (3.4,-1.8) {Hạng Nhất};
      \draw[->, thick, hustdark] (g3.east) -- node[above, font=\tiny]{120 tháng} (t3.west);
    \end{tikzpicture}
  \end{columns}
\end{frame}

% --- Slide 8: Kiến trúc tổng quan ---
\begin{frame}{Kiến trúc tổng quan hệ thống}
  \centering
  \begin{block}{Luồng dữ liệu chính}
    \centering\normalsize
    \textcolor{hustred}{\textbf{Browser}} $\xrightarrow{\text{HTTPS}}$
    \textcolor{hustred}{\textbf{Frontend}} $\xrightarrow{\text{REST}}$
    \textcolor{hustred}{\textbf{Backend}} $\xrightarrow{\text{Prisma}}$
    \textcolor{hustred}{\textbf{PostgreSQL}}\\[0.2cm]
    \footnotesize \textcolor{hustdark}{\textit{Socket.IO --- thông báo real-time hai chiều Browser $\leftrightarrow$ Backend}}
  \end{block}

  \vspace{0.3cm}
  \begin{columns}[T,onlytextwidth]
    \column{0.33\textwidth}
    \textbf{\textcolor{hustred}{Tầng giao diện}}\\
    \scriptsize \textbf{Next.js}, Ant Design, TypeScript end-to-end
    \column{0.33\textwidth}
    \textbf{\textcolor{hustred}{Tầng ứng dụng}}\\
    \scriptsize \textbf{Express} + Prisma ORM, layered\\Controller $\to$ Service $\to$ Repository
    \column{0.33\textwidth}
    \textbf{\textcolor{hustred}{Tầng dữ liệu}}\\
    \scriptsize \textbf{PostgreSQL} + Socket.IO real-time, PM2 process manager
  \end{columns}
\end{frame}

% --- Slide 9: Công nghệ sử dụng ---
\begin{frame}{Công nghệ sử dụng --- Tech Stack}
  \begin{columns}[T,onlytextwidth]
    \column{0.5\textwidth}
    \textcolor{hustred}{\textbf{Frontend}}
    \begin{itemize}\small
      \item \textbf{Next.js} (App Router) --- server components, file-based routing
      \item \textbf{TypeScript} --- type safety toàn dự án
      \item \textbf{Ant Design} --- component nghiệp vụ (Form, Table, DatePicker, Modal, ...)
      \item \textbf{Tailwind CSS} --- styling tuỳ biến cho layout, spacing
      \item \textbf{Zod} --- schema validation chia sẻ giữa client và server
      \item \textbf{Socket.IO client} --- nhận thông báo real-time
    \end{itemize}

    \column{0.5\textwidth}
    \textcolor{hustred}{\textbf{Backend \& CSDL}}
    \begin{itemize}\small
      \item \textbf{Express + TypeScript} --- REST API
      \item \textbf{Prisma ORM} --- type-safe queries, migration tự động
      \item \textbf{PostgreSQL} --- RDBMS, transaction \texttt{SERIALIZABLE}
      \item \textbf{Zod} --- validate hai tầng (BE + FE) cùng phong cách schema
      \item \textbf{JWT 2 token} --- access + refresh
      \item \textbf{bcrypt} --- hash password
    \end{itemize}
  \end{columns}

  \vspace{0.3cm}
  \centering
  \hustpill{Socket.IO} \hustpill{PM2} \hustpill{node-cron} \hustpill{ExcelJS} \hustpill{PDFKit} \hustpill{Local deploy}
\end{frame}

% --- Slide 10: Thiết kế CSDL ---
\begin{frame}{Thiết kế cơ sở dữ liệu}
  \begin{columns}[T,onlytextwidth]
    \column{0.5\textwidth}
    \textbf{Tổng quan}\\
    \begin{block}{}
      \centering {\Huge \textcolor{hustred}{\textbf{23}}}\\
      bảng chính, FK đầy đủ\\
      ID dạng CUID 25 ký tự
    \end{block}

    \textbf{Nhóm bảng}
    \begin{itemize}\small
      \item \textbf{Tổ chức} --- \texttt{QuanNhan}, \texttt{CoQuanDonVi}, \texttt{DonViTrucThuoc}, \texttt{ChucVu}
      \item \textbf{Đề xuất \& Hồ sơ} --- \texttt{BangDeXuat}, \texttt{DanhHieuHangNam}, \texttt{HoSoNienHan}, \texttt{HoSoCongHien}
      \item \textbf{Tài khoản} --- \texttt{TaiKhoan} (4 vai trò)
      \item \textbf{Phụ trợ} --- \texttt{SystemLog}, \texttt{ThongBao}, \texttt{FileQuyetDinh}, \texttt{SystemSetting}
    \end{itemize}

    \column{0.5\textwidth}
    \begin{block}{Điểm thiết kế đặc biệt}
      \scriptsize
      \texttt{FileQuyetDinh} liên kết với 8 bảng khen thưởng qua \textbf{khoá ngoại theo số quyết định} (khoá tự nhiên, không dùng ID surrogate).\\[0.2cm]
      PostgreSQL tự cascade khi đổi số quyết định, đảm bảo toàn vẹn ở tầng cơ sở dữ liệu.
    \end{block}

    \textbf{Composite uniques}
    \begin{itemize}\scriptsize
      \item \texttt{\scriptsize DanhHieuHangNam}\\\hspace*{0.3cm}\texttt{\scriptsize (quan\_nhan\_id, nam)}
      \item \texttt{\scriptsize KhenThuongHCCSVV}\\\hspace*{0.3cm}\texttt{\scriptsize (quan\_nhan\_id, danh\_hieu)}
      \item \texttt{\scriptsize ChucVu}\\\hspace*{0.3cm}\texttt{\scriptsize (co\_quan\_don\_vi\_id, ten\_chuc\_vu)}
    \end{itemize}
  \end{columns}
\end{frame}

% --- Slide 11: Use-case Diagram ---
\begin{frame}{Phân tích chức năng --- Use-case Diagram}
  \begin{columns}[T,onlytextwidth]
    \column{0.42\textwidth}
    \textbf{Phân quyền theo vai trò}\\[0.2cm]
    \scriptsize
    \begin{tabular}{|>{\raggedright\arraybackslash}p{2.5cm}|>{\raggedright\arraybackslash}p{2.6cm}|}
      \hline
      \rowcolor{hustred}\textcolor{white}{\textbf{Vai trò}} & \textcolor{white}{\textbf{Phạm vi}} \\ \hline
      \textbf{USER} & Cá nhân \\ \hline
      \textbf{MANAGER} & Đơn vị quản lý \\ \hline
      \textbf{ADMIN} & Toàn nghiệp vụ \\ \hline
      \textbf{SUPER\_ADMIN} & Hệ thống + quản trị \\ \hline
    \end{tabular}\\[0.3cm]

    \begin{block}{Hai lớp phân quyền}
      \scriptsize
      Middleware \texttt{requireRole} chặn ở route + \texttt{unitFilter} giới hạn dữ liệu theo đơn vị --- hai lớp độc lập, dù bypass được lớp này thì lớp kia vẫn chặn.
    \end{block}

    \column{0.55\textwidth}
    \textbf{Use-case chính theo vai trò}\\[0.2cm]
    \scriptsize
    \begin{tabular}{|>{\raggedright\arraybackslash}p{2cm}|>{\raggedright\arraybackslash}p{4cm}|}
      \hline
      \rowcolor{hustdark}\textcolor{white}{\textbf{Vai trò}} & \textcolor{white}{\textbf{Use-case chính}} \\ \hline
      SUPER\_ADMIN & Quản trị hệ thống, sao lưu, nhật ký toàn hệ thống \\ \hline
      ADMIN & Duyệt đề xuất, cấp số quyết định, quản lý quân nhân \\ \hline
      MANAGER & Tạo đề xuất, xem hồ sơ đơn vị, tính lại hồ sơ quân nhân \\ \hline
      USER & Xem hồ sơ cá nhân, nhận thông báo \\ \hline
    \end{tabular}
  \end{columns}
\end{frame}

% --- Slide 12: Activity Diagram quy trình duyệt ---
\begin{frame}{Thiết kế hoạt động --- Quy trình duyệt đề xuất}
  \begin{columns}[T,onlytextwidth]
    \column{0.55\textwidth}
    \textbf{Các bước chính}\\[0.2cm]
    \footnotesize
    \begin{enumerate}
      \item USER tạo đề xuất (form hoặc Excel)
      \item MANAGER duyệt cấp đơn vị
      \item ADMIN kiểm tra điều kiện và cấp số quyết định
      \item Sinh PDF, lưu file đính kèm
      \item Cập nhật hồ sơ, gửi thông báo, ghi nhật ký
    \end{enumerate}
    \vspace{0.2cm}
    \begin{block}{Tính nhất quán}
      \scriptsize Bước 3--5 thực hiện trong \textbf{một transaction Prisma} bảo đảm tính nhất quán của dữ liệu.
    \end{block}

    \column{0.42\textwidth}
    \begin{block}{Sơ đồ flow đơn giản}
      \scriptsize\centering
      \textbf{USER tạo đề xuất}\\
      $\downarrow$\\
      \textbf{MANAGER duyệt}\\
      $\downarrow$\\
      \textbf{ADMIN kiểm tra}\\
      $\downarrow$\\
      \textbf{Sinh PDF + lưu}\\
      $\downarrow$\\
      \textbf{Thông báo + nhật ký}
    \end{block}
  \end{columns}
\end{frame}

% --- Slide 13: Strategy Pattern ⭐ ---
\begin{frame}[fragile]{Kiến trúc phần mềm --- Strategy Pattern}
  \begin{columns}[T,onlytextwidth]
    \column{0.5\textwidth}
    \textbf{Vấn đề}\\
    \scriptsize 5 nhóm UC5--UC9 hiện thực thành 7 proposal type, mỗi loại có logic submit/validate/import khác nhau $\to$ \texttt{if/else} sẽ tạo file 2000+ dòng, khó test và bảo trì.\\[0.3cm]

    \textbf{Giải pháp}\\
    \scriptsize Định nghĩa \textbf{interface chung}, mỗi loại 1 file riêng, dispatch qua \textbf{Registry}.

    \begin{lstlisting}[basicstyle=\ttfamily\tiny\color{textmain}]
interface ProposalStrategy {
  buildSubmitPayload(data);
  validateApprove(ctx);
  importInTransaction(tx);
  buildSuccessMessage(result);
}
    \end{lstlisting}

    \column{0.5\textwidth}
    \begin{block}{7 Strategy implementations}
      \scriptsize\centering
      \textbf{\textcolor{hustred}{ProposalStrategy}} \textit{(interface)}\\[0.2cm]
      $\swarrow \quad \downarrow \quad \searrow$\\[0.2cm]
      \begin{tabular}{ll}
        caNhanHangNam & donViHangNam \\
        hccsvv & hcbvtq \\
        hcqkqt & knc \\
        \multicolumn{2}{c}{nckh} \\
      \end{tabular}
    \end{block}

    \begin{block}{DRY: singleMedalImporter}
      \scriptsize HCQKQT và KNC dùng chung helper để loại bỏ trùng lặp logic.
    \end{block}
  \end{columns}
\end{frame}

% --- Slide 14: Tự động xét điều kiện ---
\begin{frame}{Tự động xét điều kiện khen thưởng}
  \begin{columns}[T,onlytextwidth]
    \column{0.6\textwidth}
    \textbf{Cách thức hoạt động}\\[0.2cm]

    \begin{block}{Xét điều kiện hai lớp đồng bộ}
      \scriptsize \texttt{computeEligibilityFlags} (tính lại hồ sơ) và \texttt{checkAwardEligibility} (kiểm tra ở API) dùng chung helper lõi \texttt{chainEligibility} --- bảo đảm hiển thị và xét duyệt nhất quán.
    \end{block}

    \begin{block}{Gợi ý đề xuất tự động theo năm}
      \scriptsize Hồ sơ hằng năm tự sinh \texttt{goi\_y} như \textit{``Đủ điều kiện đề nghị BKBQP năm 2026''} --- cán bộ chỉ cần xác nhận để tạo đề xuất.
    \end{block}

    \begin{block}{Tính lại tự động khi có thay đổi}
      \scriptsize Khi quân nhân nhận danh hiệu mới, đổi đơn vị hoặc cập nhật năm phục vụ --- hệ thống tính lại toàn bộ chuỗi và cập nhật trạng thái đủ/không đủ điều kiện.
    \end{block}

    \column{0.4\textwidth}
    \begin{block}{Loại bỏ tính thủ công}
      \scriptsize \centering Loại bỏ hoàn toàn việc \textbf{tính thủ công} các quy định phức tạp 2y / 3y / 7y.
    \end{block}

    \begin{block}{Lịch sử dạng timeline}
      \scriptsize Hiển thị toàn bộ danh hiệu của một quân nhân theo trục thời gian --- dễ quan sát chuỗi BKBQP/CSTĐTQ/BKTTCP.
    \end{block}
  \end{columns}
\end{frame}

% --- Slide 15: Quy trình đề xuất + ra quyết định ---
\begin{frame}{Quy trình đề xuất và ra quyết định}
  \begin{columns}[T,onlytextwidth]
    \column{0.55\textwidth}
    \begin{block}{Form đề xuất theo từng nhóm UC}
      \scriptsize Mỗi nhóm UC5--UC9 có form riêng. Schema \textbf{Zod chia sẻ giữa client và server} đảm bảo nhất quán.
    \end{block}

    \begin{block}{Workflow ba cấp duyệt}
      \scriptsize USER $\to$ MANAGER (cấp đơn vị) $\to$ ADMIN (cấp toàn nghiệp vụ). Mỗi bước có thể duyệt hoặc từ chối kèm lý do, thông báo về người tạo.
    \end{block}

    \begin{block}{Cấp số quyết định + Sinh PDF}
      \scriptsize Sinh số quyết định theo định dạng và bộ đếm riêng từng năm --- không trùng lặp. PDF kết xuất theo template từng loại danh hiệu, lưu vào \texttt{storage/}.
    \end{block}

    \column{0.4\textwidth}
    \begin{block}{Bảo đảm toàn vẹn}
      \scriptsize \centering Toàn bộ quy trình thực hiện trong \textbf{một transaction}, kèm \textbf{audit log} chi tiết và phát thông báo real-time.
    \end{block}
  \end{columns}
\end{frame}

% --- Slide 16: Import / Export Excel ---
\begin{frame}{Nhập / Xuất dữ liệu hàng loạt qua Excel}
  \begin{columns}[T,onlytextwidth]
    \column{0.55\textwidth}
    \textbf{Pattern hai bước an toàn}\\[0.2cm]
    \begin{block}{Bước 1 --- Xem trước}
      \scriptsize Đọc file Excel, kiểm tra từng dòng (định dạng ngày, CCCD, mã quân nhân, mã đơn vị). Hiển thị bảng lỗi cụ thể trên giao diện trước khi ghi.
    \end{block}

    \begin{block}{Bước 2 --- Xác nhận}
      \scriptsize Cán bộ xác nhận dữ liệu đúng $\to$ ghi vào CSDL trong một transaction. Nếu một dòng lỗi, toàn bộ rollback.
    \end{block}

    \begin{block}{Template chuẩn có data validation}
      \scriptsize Hệ thống sinh template Excel với dropdown (cấp bậc, đơn vị, danh hiệu), định dạng ngày tháng --- cán bộ điền theo mẫu, giảm sai sót.
    \end{block}

    \column{0.4\textwidth}
    \begin{block}{Hiệu năng}
      \scriptsize Áp dụng \textbf{batch query}: collect IDs $\to$ 1 query \texttt{findMany} $\to$ đẩy vào \texttt{Map} tra cứu O(1), tránh N+1 truy vấn.
    \end{block}

    \begin{block}{Hỗ trợ}
      \scriptsize Import hàng loạt các nhóm UC5--UC8 qua Excel + Export báo cáo theo năm và đơn vị.
    \end{block}
  \end{columns}
\end{frame}

% --- Slide 17: Vận hành & Quản trị ---
\begin{frame}{Vận hành \& Quản trị hệ thống}
  \begin{columns}[T,onlytextwidth]
    \column{0.33\textwidth}
    \textbf{\textcolor{hustred}{Thông báo Real-time}}
    \begin{block}{Socket.IO push}
      \scriptsize Phát qua Socket.IO theo room theo user khi có sự kiện.
    \end{block}
    \begin{block}{Lưu lâu dài}
      \scriptsize Đồng thời lưu vào bảng \texttt{ThongBao} để người dùng offline vẫn nhận khi đăng nhập lại.
    \end{block}

    \column{0.33\textwidth}
    \textbf{\textcolor{hustred}{Audit Log}}
    \begin{block}{Ghi nhận đầy đủ}
      \scriptsize Middleware tự ghi: ai, làm gì, đối tượng nào, kèm IP \& user-agent.
    \end{block}
    \begin{block}{Lọc đa chiều}
      \scriptsize Lọc theo user, hành động, resource, khoảng thời gian.
    \end{block}

    \column{0.33\textwidth}
    \textbf{\textcolor{hustred}{Backup \& DevZone}}
    \begin{block}{Backup tự động}
      \scriptsize Lập lịch qua \texttt{node-cron}, kết xuất file SQL vào \texttt{backups/}. Restore bằng \texttt{psql}.
    \end{block}
    \begin{block}{DevZone}
      \scriptsize SUPER\_ADMIN bật/tắt schedule, trigger backup thủ công.
    \end{block}
  \end{columns}
\end{frame}

% --- Slide 18: Kiểm thử + Hiệu năng + Triển khai ⭐ ---
\begin{frame}{Kiểm thử, Hiệu năng và Triển khai}
  \begin{columns}[T,onlytextwidth]
    \column{0.55\textwidth}
    \textbf{Kiểm thử \& Hiệu năng}\\[0.2cm]

    \begin{block}{}
      \centering {\LARGE \textcolor{hustred}{\textbf{79}}} file kiểm thử Jest, pass \textbf{100\%} \quad $|$ \quad {\large \textcolor{hustred}{\textbf{N+1 $\to$ 1}}} batch query Excel
    \end{block}

    \scriptsize\textbf{4 kịch bản nổi bật:}
    \begin{enumerate}\scriptsize
      \item \textbf{Tranh chấp ghi} --- 2 admin duyệt cùng lúc, chỉ 1 thành công
      \item \textbf{Xét điều kiện BKTTCP} --- cửa sổ trượt 7 năm + ràng buộc lifetime
      \item \textbf{Chống giả mạo} --- admin sửa dữ liệu gửi lên, server vẫn ghi đúng
      \item \textbf{Vòng đời thực tế} --- quân nhân thăng tiến qua các danh hiệu
    \end{enumerate}

    \column{0.4\textwidth}
    \textbf{Triển khai mạng nội bộ}\\[0.2cm]
    \scriptsize
    \begin{enumerate}
      \item Cài \texttt{node\_modules} trên máy có Internet, chuyển sang máy nội bộ
      \item Tạo CSDL PostgreSQL rỗng + cấu hình \texttt{.env}
      \item \texttt{npm run setup} --- tạo bảng và SUPER\_ADMIN (idempotent)
      \item \texttt{npm run serve} --- build + khởi chạy bằng PM2
    \end{enumerate}

    \begin{block}{PM2 production}
      \scriptsize Auto-restart, RAM 500MB, log rotation. Hệ thống vận hành \textbf{hoàn toàn offline} sau khi cài đặt.
    \end{block}
  \end{columns}
\end{frame}

% --- Slide 19: Kết luận và Hướng phát triển ---
\begin{frame}{Kết luận và hướng phát triển}
  \begin{columns}[T,onlytextwidth]
    \column{0.33\textwidth}
    \textbf{\textcolor{hustred}{Đã đạt được}}
    \begin{itemize}\scriptsize
      \item Hoàn thành 5 nhóm khen thưởng UC5--UC9, 4 vai trò
      \item Tự động hoá quy trình xét duyệt
      \item 79 file test, pass 100\%
      \item Triển khai trên mạng nội bộ
      \item Layered + Strategy pattern
    \end{itemize}

    \column{0.33\textwidth}
    \textbf{\textcolor{hustred}{Hạn chế}}
    \begin{itemize}\scriptsize
      \item BKTTCP cá nhân lifetime --- chưa hỗ trợ danh hiệu cao hơn
      \item Chưa có module BI thống kê chi tiết
      \item Chỉ chạy 1 instance --- chưa scale ngang
      \item Ký số còn dưới dạng file PDF
    \end{itemize}

    \column{0.33\textwidth}
    \textbf{\textcolor{hustred}{Hướng phát triển}}
    \begin{itemize}\scriptsize
      \item Ứng dụng di động cho cấp duyệt nhanh
      \item Module BI: biểu đồ xu hướng theo năm \& đơn vị
      \item Tích hợp \textbf{Smart Card / USB token}
      \item Cluster mode + Redis adapter Socket.IO
    \end{itemize}
  \end{columns}

  \vspace{0.3cm}
  \begin{block}{}
    \centering\small Phần mềm đã sẵn sàng \highlight{triển khai thử nghiệm tại đơn vị}, đáp ứng đầy đủ các mục tiêu đặt ra ban đầu.
  \end{block}
\end{frame}

% --- Slide 20: Cảm ơn ---
{
\setbeamertemplate{footline}{}
\begin{frame}[plain]
  \begin{center}
    \vspace{1.5cm}
    {\huge \textcolor{hustred}{\textbf{EM XIN TRÂN TRỌNG CẢM ƠN}}}\\[1cm]

    {\Large Hội đồng đã lắng nghe phần trình bày}\\[1.5cm]

    {\large Em rất mong nhận được \highlight{nhận xét và câu hỏi} từ Hội đồng}\\[2cm]

    {\Huge \textcolor{hustred}{$\heartsuit$}}\\
    {\Large \textcolor{hustdark}{Q\&A}}
  \end{center}
\end{frame}
}

\end{document}
